\documentclass[10pt,a4paper]{amsart}
\usepackage[margin=19mm]{geometry}
\usepackage{amsmath,amssymb,mathtools}
\usepackage[scr=rsfs,cal=euler]{mathalfa}
\usepackage{xfrac}
\usepackage[unicode,hidelinks,bookmarks=false]{hyperref}
\newcommand{\ie}{i.e.\ }
\newcommand{\cf}{cf.\ }
\newcommand{\resp}{respectively\ }
\newcommand{\Subsection}{\subsection}
\providecommand{\nobreakdash}{\nobreak\hbox{-}}
\setlength{\emergencystretch}{2em}
\multlinegap0pt
\usepackage{lmodern}
\newtheorem{theorem}{Theorem}
\newtheorem{lemma}[theorem]{Lemma}
\newtheorem{corollary}[theorem]{Corollary}
\newtheorem{remark}[theorem]{Remark}
\newcommand{\E}{\mathbb E}
\newcommand{\Pp}{\mathbb P}
\newcommand{\Bin}{\operatorname{Bin}}
\newcommand{\Fint}{\mathcal F^{\circ}}
\begin{document}
\title[Reflection of optimal supports for k-wise independent bits]{Reflection of optimal supports for k-wise independent bits}
\author{Roy Hermann}
\date{}
\maketitle
\noindent\emph{Source and licence.} Reflection of optimal supports for k-wise independent bits. Version arXiv:2609.25595v1 (CC BY 4.0). This material is distributed under the Creative Commons Attribution 4.0 International licence, http://creativecommons.org/licenses/by/4.0/, as recorded in the arXiv OAI licence metadata for this exact version and shown on the abstract page https://arxiv.org/abs/2609.25595v1. Full rights notice: licenses/arxiv-2609.25595-CC-BY-4.0.txt.\par\smallskip
\noindent\textit{Editorial scope.} Counted unit: the original \texttt{theorem} environment \texttt{thm:identity} at source lines 57--69 together with its complete proof at lines 71--112; the delivered document keeps source lines 42--113 so that every referenced label and every citation resolves inside it. Native editable source is the paper's own concatenated TeX; the AMS wrapper is editorial formatting only. No publisher class, upstream script or unrelated artwork is redistributed.
\begin{equation}\label{eq:basis}
 I=J\cup\{n\},\qquad
 J=\bigcup_{j=1}^{k/2}\{a_j,a_j+1\}\subseteq\{0,\ldots,n-1\},
\end{equation}
where $a_{j+1}\ge a_j+2$. For such an $I$, denote by $w_s^I(p)$ the unique signed solution of the moment equalities supported on $I$.
Uniqueness follows because the polynomials $\binom{x}{r}$, $0\le r\le k$, form a basis for polynomials of degree at most $k$, and evaluation on $k+1$ distinct nodes is invertible. We call $I$ feasible at $p$ if all its basic weights are nonnegative.

Define the involution
\begin{equation}\label{eq:star}
 J^*=\{n-1-s:s\in J\},\qquad I^*=J^*\cup\{n\}.
\end{equation}
It takes each pair $\{a,a+1\}$ to $\{n-a-2,n-a-1\}$ and preserves the class~\eqref{eq:basis}. The node $n$ is fixed separately.

\section{Reflection of the basic weights}


\begin{theorem}\label{thm:identity}
For every basis~\eqref{eq:basis}, every $s\in J$, and every $p\in(0,1)$,
\begin{equation}\label{eq:identity}
 w_{n-1-s}^{I^*}(1-p)
 =\frac{p(n-s)}{(1-p)(s+1)}\,w_s^I(p).
\end{equation}
Furthermore, $w_n^I(p)\ge p^n>0$.
In particular, if
\[
 \Fint(I)=\{p\in(0,1):w_s^I(p)\ge0\text{ for all }s\in I\},
\]
then $\Fint(I^*)=\{1-p:p\in\Fint(I)\}$.
\end{theorem}

\begin{proof}
For $s\in J$, let
\[
 \ell_s^J(x)=\prod_{t\in J\setminus\{s\}}\frac{x-t}{s-t}
\]
be the Lagrange polynomial on $J$. The Lagrange polynomial for the same node on $I$ is
\begin{equation}\label{eq:lagrange}
 L_s^I(x)=\frac{n-x}{n-s}\ell_s^J(x).
\end{equation}
Let $X\sim\Bin(n,p)$. The moment equalities imply exact agreement with the binomial law on every polynomial of degree at most $k$. Since $L_s^I$ has that degree and is one at $s$ and zero at the other nodes of $I$,
\begin{equation}\label{eq:expectation}
 w_s^I(p)=\E L_s^I(X).
\end{equation}
For any function $f$ on $\{0,\ldots,n\}$, the identity
$(n-x)\binom{n}{x}=n\binom{n-1}{x}$ gives
\begin{equation}\label{eq:reweight}
 \E[(n-X)f(X)]=n(1-p)\E f(Y),\qquad Y\sim\Bin(n-1,p).
\end{equation}
Consequently,
\begin{equation}\label{eq:reduced}
 w_s^I(p)=\frac{n(1-p)}{n-s}\E\ell_s^J(Y).
\end{equation}
Writing $s^*=n-1-s$, interpolation gives
\[
 \ell_{s^*}^{J^*}(x)=\ell_s^J(n-1-x).
\]
If $Y^*\sim\Bin(n-1,1-p)$, then $n-1-Y^*$ has the same law as $Y$. Applying~\eqref{eq:reduced} to $I^*$ at $1-p$ yields
\[
 w_{s^*}^{I^*}(1-p)
 =\frac{np}{s+1}\E\ell_{s^*}^{J^*}(Y^*)
 =\frac{np}{s+1}\E\ell_s^J(Y).
\]
Comparison with~\eqref{eq:reduced} proves~\eqref{eq:identity}, including at zeros of the weights.

The Lagrange polynomial for the distinguished node is
\begin{equation}\label{eq:dual}
 Q_I(x)=\prod_{j\in J}\frac{x-j}{n-j},\qquad
 w_n^I(p)=\E Q_I(X).
\end{equation}
For an integer $x$, every paired factor $(x-a_j)(x-a_j-1)$ is nonnegative, and every denominator $n-j$ is positive. Hence $Q_I(x)\ge0$ for $x\in\{0,\ldots,n\}$, with $Q_I(n)=1$. This proves $w_n^I(p)\ge p^n>0$.
All the multipliers in~\eqref{eq:identity} are positive on $(0,1)$, so the corresponding other weights have the same signs. Applying the same argument to $I^*$ proves the asserted equality of feasibility sets.
\end{proof}

\end{document}
